\documentclass[11pt]{article}
\usepackage[margin=1in]{geometry}
\usepackage{amsmath,amssymb,amsthm}
\usepackage{algorithm}
\usepackage{algpseudocode}
\usepackage{booktabs}
\usepackage{hyperref}

\newtheorem{definition}{Definition}
\newtheorem{theorem}{Theorem}
\newtheorem{example}{Example}

\title{Interval Belief Propagation\\for Credal Networks with Multi-Valued Variables}
\author{IBM Research}
\date{}

\begin{document}
\maketitle

\begin{abstract}
Interval Belief Propagation (IBP) is an iterative message-passing algorithm for
approximate marginal inference in credal networks. Unlike the 2U and L2U algorithms
which are restricted to binary variables, IBP operates natively on multi-valued
variables by propagating interval-vector messages. At each iteration, variable nodes
tighten their interval bounds by intersecting incoming messages, while factor nodes
compute new bounds by optimizing over the extreme points of local credal sets and the
corner points of incoming interval messages. The algorithm is inspired by the ApproxLP
framework of Antonucci et al.\ (2015) and the modified belief propagation for Logical
Credal Networks (Qian et al., 2021).
\end{abstract}

\section{Preliminaries}

\subsection{Factor Graph}

Given a credal network $\langle G, \mathcal{K} \rangle$, we construct a
\textbf{factor graph}: a bipartite graph with variable nodes and factor nodes.
Each variable $X_i$ in the credal network corresponds to a variable node with
cardinality $|D_i|$. Each conditional credal set $K(X_i \mid \mathrm{Pa}(X_i))$
defines a factor node $f_i$ whose scope is $\{X_i\} \cup \mathrm{Pa}(X_i)$.
The factor stores the extreme points of $K(X_i \mid \pi)$ for all parent
configurations $\pi$.

\subsection{Interval Messages}

\begin{definition}[Interval Message]
For a variable $X$ with $k$ states, an \textbf{interval message} is a pair of
vectors $(\mathbf{l}, \mathbf{u}) \in \mathbb{R}^k \times \mathbb{R}^k$ where
$0 \leq l_x \leq u_x \leq 1$ for each state $x \in D_X$. The interval
$[l_x, u_x]$ bounds the probability $P(X = x)$ from the message sender's
perspective.
\end{definition}

Messages are of two types:
\begin{itemize}
  \item $\mu_{v \to f} = (\mathbf{l}_{v \to f}, \mathbf{u}_{v \to f})$:
        message from variable $v$ to factor $f$.
  \item $\mu_{f \to v} = (\mathbf{l}_{f \to v}, \mathbf{u}_{f \to v})$:
        message from factor $f$ to variable $v$.
\end{itemize}

All messages are initialized to the vacuous interval $([0, \dots, 0], [1, \dots, 1])$.

\section{Algorithm}

\begin{algorithm}[H]
\caption{Interval Belief Propagation (IBP)}
\label{alg:ibp}
\begin{algorithmic}[1]
\Require Factor graph with extreme points, query $Z$, evidence $\mathbf{Y} = \mathbf{y}$,
         max iterations $T$, threshold $\tau$
\Ensure Bounds $[\underline{P}(Z{=}z),\; \overline{P}(Z{=}z)]$ for each $z$
\Statex
\State Initialize $\mu_{v \to f} \gets (\mathbf{0}, \mathbf{1})$ and
       $\mu_{f \to v} \gets (\mathbf{0}, \mathbf{1})$ for all edges
\For{each evidence variable $Y_j = y_j$}
  \State Set $\mu_{Y_j \to f} \gets (\delta_{y_j}, \delta_{y_j})$ for all factors
         $f \in \mathcal{N}(Y_j)$
\EndFor
\Statex
\For{$t = 1, \dots, T$}
  \State $\delta \gets 0$
  \Statex
  \Comment{\textsc{Variable-to-Factor Updates}}
  \For{each variable $v$ (not evidence), each factor $f \in \mathcal{N}(v)$}
    \For{each state $x$}
      \State $l_{v \to f}(x) \gets \max_{f' \in \mathcal{N}(v) \setminus \{f\}}
             l_{f' \to v}(x)$
      \State $u_{v \to f}(x) \gets \min_{f' \in \mathcal{N}(v) \setminus \{f\}}
             u_{f' \to v}(x)$
    \EndFor
    \State $\mathbf{u}_{v \to f} \gets \max(\mathbf{l}_{v \to f}, \mathbf{u}_{v \to f})$
    \Comment{Ensure $l \leq u$}
    \State Update $\delta$
  \EndFor
  \Statex
  \Comment{\textsc{Factor-to-Variable Updates}}
  \For{each factor $f_i$ (node $X_i$, parents $\mathrm{Pa}(X_i)$), each $v \in \mathrm{scope}(f_i)$}
    \State $\mathbf{l}_{\mathrm{new}} \gets \mathbf{1}, \quad \mathbf{u}_{\mathrm{new}} \gets \mathbf{0}$
    \For{each vertex combination $\mathbf{vc}$ of $\mathrm{ext}\,K(X_i \mid \cdot)$}
      \For{each corner-point combination $\mathbf{cc}$ of incoming intervals from
           $\mathrm{scope}(f_i) \setminus \{v\}$}
        \State Compute unnormalized marginal $\mathbf{m}(v)$ using $\mathbf{vc}$ and $\mathbf{cc}$
        \Comment{See Section~\ref{sec:ftov}}
        \State $\mathbf{p} \gets \mathbf{m} / \sum_{x} m(x)$
        \State $\mathbf{l}_{\mathrm{new}} \gets \min(\mathbf{l}_{\mathrm{new}}, \mathbf{p})$
        \State $\mathbf{u}_{\mathrm{new}} \gets \max(\mathbf{u}_{\mathrm{new}}, \mathbf{p})$
      \EndFor
    \EndFor
    \State $\mu_{f \to v} \gets (\mathbf{l}_{\mathrm{new}}, \mathbf{u}_{\mathrm{new}})$
    \State Update $\delta$
  \EndFor
  \If{$\delta < \tau$} \textbf{break} \EndIf
\EndFor
\Statex
\Comment{\textsc{Extract Marginal Bounds}}
\For{each state $z$}
  \State $\underline{P}(z) \gets \max_{f \in \mathcal{N}(Z)} l_{f \to Z}(z)$
  \State $\overline{P}(z) \gets \min_{f \in \mathcal{N}(Z)} u_{f \to Z}(z)$
\EndFor
\State \Return $[\underline{P}(z), \overline{P}(z)]$ for each $z$
\end{algorithmic}
\end{algorithm}

\subsection{Factor-to-Variable Message Computation}\label{sec:ftov}

The factor-to-variable message from factor $f_i$ (associated with node $X_i$,
parents $\mathrm{Pa}(X_i)$) to target variable $v \in \mathrm{scope}(f_i)$
is computed as follows.

Let $\mathcal{O} = \mathrm{scope}(f_i) \setminus \{v\}$ be the other variables.
For each ``other'' variable $w \in \mathcal{O}$, let
$(\mathbf{l}_{w}, \mathbf{u}_{w}) = \mu_{w \to f_i}$ be the incoming message.

\paragraph{Vertex enumeration.}
The factor stores extreme point vertices $V(\pi) = \mathrm{ext}\,K(X_i \mid \pi)$
for each parent configuration $\pi$. A vertex combination selects one vertex per
parent configuration.

\paragraph{Corner-point enumeration.}
For each variable $w \in \mathcal{O}$, the incoming interval $[l_w(x), u_w(x)]$
is evaluated at its two extreme values: $l_w(x)$ or $u_w(x)$. With $|\mathcal{O}|$
variables, there are $2^{|\mathcal{O}|}$ corner-point combinations.

\paragraph{Marginal computation.}
For a given vertex combination $\mathbf{vc}$ and corner-point assignment
$\mathbf{cc}$, the unnormalized marginal over $v$ is:

\emph{Case 1: $v = X_i$ (target is the child node).}
\[
  m(v{=}x) = \sum_{\pi} P_{\mathbf{vc}}(X_i{=}x \mid \pi)
  \cdot \prod_{w \in \mathrm{Pa}(X_i) \cap \mathcal{O}} q_w(\pi_w),
\]
where $q_w$ is the lower or upper bound vector selected by $\mathbf{cc}$ for
variable $w$, and $\pi_w$ is the value of $w$ in configuration $\pi$.

\emph{Case 2: $v \in \mathrm{Pa}(X_i)$ (target is a parent).}
\[
  m(v{=}x) = \sum_{\pi : \pi_v = x}
  \Bigl(\sum_{x_i} P_{\mathbf{vc}}(X_i{=}x_i \mid \pi) \cdot q_{X_i}(x_i)\Bigr)
  \cdot \prod_{w \in \mathrm{Pa} \setminus \{v\}} q_w(\pi_w),
\]
where $q_{X_i}$ is the incoming child message bound if $X_i \in \mathcal{O}$.

The marginal is then normalized: $p(x) = m(x) / \sum_{x'} m(x')$.

\section{Convergence Properties}

\begin{itemize}
  \item Messages are updated by taking min/max over all vertex and corner-point
        combinations, so the algorithm explores the full range of the credal set.
  \item On polytree-shaped networks with binary variables, IBP reduces to the
        2U algorithm and converges to exact bounds.
  \item On loopy graphs, convergence to a fixed point is empirically observed
        but not formally guaranteed. The fixed point may not correspond to the
        exact bounds (similar to loopy belief propagation for Bayesian networks).
\end{itemize}

\section{Running Example}

\begin{example}[Alarm Network]
Using the same alarm network with variables $B$ (2 states), $E$ (2 states),
$A$ (2 states), and $CD$ (4 states).

\paragraph{Factor graph.}
Four factors: $f_A$ (scope $\{A, B, E\}$), $f_B$ (scope $\{B\}$),
$f_E$ (scope $\{E\}$), $f_{CD}$ (scope $\{CD, A\}$).

\paragraph{Query: $P(B{=}1)$, no evidence.}
Variable $B$ connects to factors $f_B$ and $f_A$.

\textbf{Iteration 0:}
All messages are vacuous $([0,0], [1,1])$.

$f_B \to B$: Factor $f_B$ has scope $\{B\}$ with vertices $[0.8, 0.2]$ and $[0.9, 0.1]$.
No other variables in scope, so: $l_{f_B \to B} = [0.8, 0.1]$, $u_{f_B \to B} = [0.9, 0.2]$.
These are the correct marginal bounds for $B$ from its local credal set.

$f_A \to B$: Factor $f_A$ has scope $\{A, B, E\}$. With vacuous messages from $A$ and $E$,
the message to $B$ is computed by enumerating vertex and corner-point combinations.
Since $A$ and $E$ messages are vacuous $[0,1]$, the resulting bounds for $B$ are wide.

\textbf{Final marginal for $B$:}
$\underline{P}(B{=}1) = \max(l_{f_B \to B}(1), l_{f_A \to B}(1))$,
$\overline{P}(B{=}1) = \min(u_{f_B \to B}(1), u_{f_A \to B}(1))$.
The bounds from $f_B$ provide the tightest constraints: $[0.1, 0.2]$.
The bounds from $f_A$ are wider and do not further constrain the result.

\paragraph{Comparison with exact VE.}
\begin{center}
\begin{tabular}{lcc}
\toprule
\textbf{Method} & $\underline{P}(B{=}1)$ & $\overline{P}(B{=}1)$ \\
\midrule
Exact VE & 0.1000 & 0.2000 \\
Interval BP & 0.1000 & 0.2000 \\
\bottomrule
\end{tabular}
\end{center}

For root nodes, IBP recovers the exact bounds because the local factor alone
determines the marginal. For non-root nodes with evidence, the bounds may differ
due to the approximation inherent in loopy message passing.
\end{example}

\section{Multi-Valued Variable Support}

A key advantage of IBP over L2U and GL2U is native support for multi-valued
variables:

\begin{itemize}
  \item \textbf{L2U}: Restricted to binary variables. Messages are scalar intervals.
  \item \textbf{GL2U}: Handles multi-valued variables by binarization (replacing a
        $k$-valued variable with $\lceil \log_2 k \rceil$ binary variables in a chain).
        This introduces auxiliary structure and can be lossy.
  \item \textbf{IBP}: Messages are interval vectors $(\mathbf{l}, \mathbf{u}) \in
        \mathbb{R}^k \times \mathbb{R}^k$. No binarization needed. Each state is
        bounded independently.
\end{itemize}

For example, the compound variable $CD$ in the alarm network has 4 states. IBP
directly propagates 4-dimensional interval vectors for $CD$, while GL2U would
require binarizing $CD$ into two binary variables.

\section{Complexity}

Each iteration of IBP takes $O(N_f \cdot V_{\max} \cdot 2^{S_{\max}} \cdot k^{S_{\max}})$
time, where $N_f$ is the number of factors, $V_{\max}$ is the maximum number of vertex
combinations per factor, $S_{\max}$ is the maximum scope size, and $k$ is the maximum
domain size. The number of iterations until convergence depends on the network structure
and is typically 10--100 for practical networks.

\section{References}

\begin{enumerate}
\item A.\ Antonucci, C.P.\ de Campos, D.\ Huber, M.\ Zaffalon. Approximate credal
      network updating by linear programming with applications to decision making.
      \emph{IJAR}, 58:25--38, 2015.
\item J.S.\ Ide, F.G.\ Cozman. IPE and L2U: Approximate algorithms for credal
      networks. \emph{STAIRS}, 2004.
\item A.\ Antonucci, Y.\ Sun, C.P.\ de Campos, M.\ Zaffalon. Generalized loopy 2U:
      A new algorithm for approximate inference in credal networks. \emph{IJAR},
      51(5):474--484, 2010.
\item H.\ Qian, R.\ Marinescu et al. Logical Credal Networks. \emph{AAAI}, 2021.
\item D.D.\ Mau\'a, F.G.\ Cozman. Thirty years of credal networks. \emph{IJAR},
      126:133--157, 2020.
\end{enumerate}

\end{document}
